\abstractPT  % Do NOT modify this line

O objetivo deste trabalho é desenvolver uma aplicação que execute operações sobre uma base de dados relacional fornecida pelos docentes. Para além da base de dados inicial, foi necessário implementar funcionalidades adicionais em SQL, nomeadamente triggers, uma view com suporte a operações de inserção e atualização, e um procedimento armazenado.
A aplicação foi implementada em Java, recorrendo à API Jakarta Persistence (JPA) para a gestão da persistência.
A JPA é uma API padronizada da plataforma Jakarta EE que facilita o mapeamento entre objetos Java e tabelas relacionais. Permite realizar operações de persistência (inserção, atualização, remoção, consulta) através de um modelo orientado a objetos, abstraindo o uso direto de SQL. Além disso, fornece suporte a transações, consultas com JPQL e controlo de concorrência através de mecanismos como locking otimista.
